% ============================================================================
% MUSTERLÖSUNG - Sequenz 1a: Rally + Begrüßung
% ============================================================================

\documentclass[11pt,a4paper]{article}

\usepackage[utf8]{inputenc}
\usepackage[T1]{fontenc}
\usepackage[ngerman]{babel}
\usepackage[left=2cm,right=2cm,top=1.8cm,bottom=1.5cm]{geometry}
\usepackage{helvet}
\usepackage{xcolor}
\usepackage{tabularx}
\usepackage{array}
\usepackage{enumitem}
\usepackage{tcolorbox}
\usepackage{fancyhdr}
\usepackage{lastpage}

\renewcommand{\familydefault}{\sfdefault}

\definecolor{spanisch}{HTML}{c41e3a}
\definecolor{grau}{HTML}{666666}
\definecolor{loesung}{HTML}{c62828}

\pagestyle{fancy}
\fancyhf{}
\fancyhead[L]{\small\textcolor{grau}{Spanisch Klasse 7}}
\fancyhead[C]{\small\textcolor{loesung}{\textbf{MUSTERLÖSUNG}}}
\fancyhead[R]{\small\textcolor{grau}{ML-01a}}
\fancyfoot[C]{\small\thepage/\pageref{LastPage}}
\renewcommand{\headrulewidth}{0.5pt}

\newcommand{\luecke}[1]{\underline{\hspace{#1}}}
\newcommand{\lsg}[1]{\textcolor{loesung}{\textbf{#1}}}
\newcommand{\punktebox}[1]{\hfill\fbox{\small \_/#1}}

\newtcolorbox{merkkasten}[1][]{
    colback=white,
    colframe=spanisch,
    fonttitle=\bfseries\color{spanisch},
    title=#1,
    boxrule=1.5pt,
    arc=0pt,
    left=5pt, right=5pt, top=5pt, bottom=5pt
}

\newcounter{aufgabe}
\newcommand{\aufgabe}[2]{%
    \stepcounter{aufgabe}%
    \vspace{0.5em}%
    \noindent\textbf{\theaufgabe.} #1 \punktebox{#2}%
    \vspace{0.3em}%
}

\begin{document}

% ============================================================================
% KOPFBEREICH
% ============================================================================
\noindent
\begin{tabularx}{\textwidth}{@{}Xr@{}}
    {\Large\textbf{\textcolor{spanisch}{¡Empezamos!}}} & Name: \luecke{5cm} \\[0.3em]
    {\small Rally + Begrüßung} & Datum: \luecke{3cm} \\
\end{tabularx}

\vspace{0.5em}
\hrule
\vspace{1em}

% ============================================================================
% MERKKASTEN
% ============================================================================
\begin{merkkasten}[Kasten 1: Begrüßung \& Verabschiedung]
\vspace{0.3em}
\begin{tabularx}{\textwidth}{@{}X|X@{}}
\textbf{Begrüßung (Saludar)} & \textbf{Verabschiedung (Despedirse)} \\[0.5em]
\hline
\lsg{¡Hola!} (Hallo) & \lsg{¡Adiós!} (Tschüss) \\[0.8em]
\lsg{Buenos días} (Guten Morgen) & ¡Hasta luego! (\lsg{Bis später}) \\[0.8em]
Buenas tardes (\lsg{Guten Tag}) & ¡Hasta mañana! (\lsg{Bis morgen}) \\[0.8em]
\lsg{Buenas noches} (Guten Abend) & \\
\end{tabularx}
\vspace{0.3em}
\end{merkkasten}

\vspace{1.5em}

% ============================================================================
% AUFGABEN MIT LÖSUNGEN
% ============================================================================

\aufgabe{Ordne zu. Verbinde die passenden Paare.}{4}

\vspace{0.5em}
\begin{center}
\begin{tabular}{ccc}
\begin{tabular}{|l|}
\hline
1. ¡Hola! \\
\hline
2. Buenos días \\
\hline
3. ¡Hasta mañana! \\
\hline
4. Buenas noches \\
\hline
\end{tabular}
&
\hspace{2cm}
&
\begin{tabular}{|l|}
\hline
a) Guten Abend \\
\hline
b) Bis morgen! \\
\hline
c) Guten Morgen \\
\hline
d) Hallo! \\
\hline
\end{tabular}
\end{tabular}
\end{center}

\vspace{0.3em}
\noindent Lösung: 1 -- \lsg{d}, \quad 2 -- \lsg{c}, \quad 3 -- \lsg{b}, \quad 4 -- \lsg{a}

\vspace{1.5em}

\aufgabe{Welche Begrüßung passt? Schreibe das richtige Wort.}{4}

\vspace{0.5em}
\begin{enumerate}[label=\alph*), nosep, leftmargin=*]
    \item 8:00 Uhr morgens: \lsg{Buenos días}
    \item 15:00 Uhr nachmittags: \lsg{Buenas tardes}
    \item 21:00 Uhr abends: \lsg{Buenas noches}
    \item Zu jeder Zeit (informell): \lsg{¡Hola!}
\end{enumerate}

\vspace{1.5em}

\aufgabe{Rally durchs Buch: Beantworte die Fragen.}{6}

\vspace{0.5em}
\begin{enumerate}[label=\alph*), nosep, leftmargin=*]
    \item Wie viele Unidades (Lektionen) hat das Buch? \lsg{8}
    \item Auf welcher Seite beginnt das Wörterverzeichnis? S. \lsg{196}
    \item Was bedeutet das Symbol \textbf{CD}? \lsg{Hörübung / Audio}
    \item Nenne 3 spanische Wörter, die du schon verstehst (Internationalismen):

    \lsg{hotel}, \quad \lsg{restaurante}, \quad \lsg{música} \textit{(o.ä.)}
\end{enumerate}

\vspace{1.5em}

\aufgabe{Schreibe einen Mini-Dialog. A begrüßt B, B antwortet.}{6}

\vspace{0.5em}
\noindent\fbox{\parbox{\textwidth-2\fboxsep-2\fboxrule}{%
\textcolor{loesung}{\textit{Beispiellösung:}}\\[0.3em]
\textbf{A:} \lsg{¡Hola! Buenos días.}\\[0.5em]
\textbf{B:} \lsg{¡Hola! Buenos días.}\\[0.5em]
\textbf{A:} \lsg{¡Hasta luego!}\\[0.5em]
\textbf{B:} \lsg{¡Adiós! ¡Hasta mañana!}
}}

\vspace{1.5em}
\hrule
\vspace{0.5em}

\begin{flushright}
\textbf{Punkte:} \luecke{1cm} / 20
\end{flushright}

\end{document}
